\documentclass[11pt]{article}

% Set 4:5 Aspect Ratio (Standard LinkedIn Portrait Carousel 1080x1350 equivalent)
\usepackage[paperwidth=8in, paperheight=10in, top=1.2in, bottom=1.2in, left=1in, right=1in]{geometry}

\usepackage{fontspec}
\usepackage[english, bidi=bidi, provide=*]{babel}

\babelprovide[import, onchar=ids fonts]{english}

% Set premium typography using Noto Sans system fonts
\babelfont{rm}{Noto Sans}
\babelfont{sf}{Noto Sans}
\babelfont{tt}{Noto Sans Mono}

% --- PACKAGES FOR MODERN DESIGN & MATHS ---
\usepackage{amsmath}
\usepackage{amssymb}
\usepackage{xcolor}
\usepackage{listings}
\usepackage{mdframed}
\usepackage{tikz}
\usetikzlibrary{shapes.geometric, arrows.meta, positioning, shadows.blur}
\usepackage{tcolorbox}
\tcbuselibrary{listings, skins, breakable}
\usepackage{fancyhdr}
\usepackage{enumitem}

% --- DESIGN TOKENS (Stripe & Linear Inspired) ---
\definecolor{brandnavy}{HTML}{0A2540}    % Deep Primary Navy
\definecolor{brandgray}{HTML}{425466}    % Slate Gray
\definecolor{brandlight}{HTML}{F8F9FA}   % Light Soft Gray
\definecolor{brandblue}{HTML}{635BFF}    % Accent Classic Blue
\definecolor{brandgreen}{HTML}{00D4B2}   % Success/Muted Green
\definecolor{brandorange}{HTML}{F25C05}  % Accent Orange
\definecolor{brandred}{HTML}{E5484D}     % Warning Red

% Premium Dark Theme for Code Blocks
\definecolor{codebackground}{HTML}{141419}
\definecolor{codekeyword}{HTML}{FF7B72}
\definecolor{codestring}{HTML}{A5D6FF}
\definecolor{codecomment}{HTML}{8B949E}
\definecolor{codetext}{HTML}{E6EDF3}
\definecolor{codeframe}{HTML}{30363D}

% --- ZERO PARAGRAPH INDENT ---
\setlength{\parindent}{0pt}
\setlength{\parskip}{1.2em}

% Custom TColorBox for styled content blocks
\newtcolorbox{infobox}[1][]{
    colback=brandlight, colframe=brandblue, arc=5pt, boxrule=0.8pt,
    left=15pt, right=15pt, top=12pt, bottom=12pt, #1
}

\newtcolorbox{darkbox}[1][]{
    colback=brandnavy, colframe=brandnavy, arc=5pt, boxrule=0.8pt,
    left=15pt, right=15pt, top=12pt, bottom=12pt, 
    coltext=white, #1
}

% Custom TColorBox for Syntax Highlighted Code
\newtcblisting{codebox}[2][]{
    colback=codebackground, colframe=codeframe, arc=5pt, boxrule=1pt,
    left=10pt, right=10pt, top=10pt, bottom=10pt,
    listing only, listing options={
        basicstyle=\ttfamily\footnotesize\color{codetext},
        keywordstyle=\color{codekeyword}\bfseries,
        stringstyle=\color{codestring},
        commentstyle=\color{codecomment}\itshape,
        showstringspaces=false,
        breaklines=true,
        language=#2
    },
    #1
}

% --- HEADER & FOOTER CONFIGURATION ---
\pagestyle{fancy}
\fancyhf{}
\renewcommand{\headrulewidth}{0pt}
\renewcommand{\footrulewidth}{0.5pt}

\fancyhead[L]{\fontsize{9}{11}\selectfont\color{brandgray}\MakeUppercase{MLOps Deep Dive}}
\fancyhead[R]{\fontsize{9}{11}\selectfont\color{brandblue}\textbf{VALIDATION DRIFT}}

\fancyfoot[L]{\fontsize{8}{10}\selectfont\color{brandgray}
Author: Dibayendu Mukherjee \hspace{0.7em}$\cdot$\hspace{0.7em}
\mbox{AI Research Engineer \hspace{0.7em}\textbullet\hspace{0.7em} Systems Engineer}}

\fancyfoot[R]{\fontsize{8}{10}\selectfont\color{brandgray}Slide \thepage}

\begin{document}

% ==========================================================
% SLIDE 1: COVER
% ==========================================================
\thispagestyle{empty}
\vspace*{0.3in}

\begin{center}

    {\fontsize{14}{16}\selectfont\color{brandblue}\textbf{
    EVALUATION \& VALIDATION
    }}

    \vspace{0.5in}

    {\fontsize{40}{48}\selectfont\color{brandnavy}\bfseries
Why your training loss\\[0.15in]
looks great but your\\[0.25in]
model still fails\\[0.15in]
in production.
}

    \vspace{0.6in}

    {\fontsize{18}{24}\selectfont\color{brandgray}
    The silent gap between validation\\
    metrics and real-world performance.
    }

    \vfill

    \begin{minipage}{0.8\textwidth}
        \centering
        \rule{0.6\linewidth}{0.5pt}

        \vspace{0.2in}

        {\fontsize{12}{14}\selectfont\color{brandnavy}
        \textbf{Dibayendu Mukherjee}
        }

        \vspace{0.05in}

        {\fontsize{10}{12}\selectfont\color{brandgray}
        AI Research Engineer \quad $\cdot$ \quad Systems Engineer
        }

    \end{minipage}

\end{center}

\newpage

% ==========================================================
% SLIDE 2: THE ILLUSION OF SUCCESS
% ==========================================================
\vspace*{0.4in}

\begin{center}
{\fontsize{26}{30}\selectfont\color{brandnavy}\bfseries 
The Illusion of Success}
\end{center}

\vspace{0.3in}

{\large\color{brandgray}
Everything looked perfect on paper.
}

\vspace{0.15in}

\begin{darkbox}
{\large\textbf{What I Saw}}

\vspace{0.05in}
\begin{itemize}[leftmargin=*, itemsep=0.08in, label={\color{brandblue}\textbullet}]
    \item 98\% validation accuracy.
    \item Smoothly decreasing training and validation loss.
    \item No signs of overfitting.
    \item Model deployed with confidence.
\end{itemize}
\end{darkbox}

\vspace{0.2in}

In production, the model failed on the most important users and edge cases. The metrics I trusted were not measuring what actually mattered.

\vspace{0.15in}

\begin{infobox}
\textbf{The Fatal Assumption:}\\
I assumed a high validation score guaranteed real-world performance. It didn't.
\end{infobox}

\vfill
\newpage

% ==========================================================
% SLIDE 3: THE MISTAKE
% ==========================================================
\section*{The Mistake}

\vspace{0.15in}

My biggest error was treating validation as a box-ticking exercise instead of a rigorous simulation of production.

\vspace{0.2in}

\begin{minipage}[t]{0.45\textwidth}
    {\large\color{brandred}\textbf{What I Did}}

    \begin{itemize}[leftmargin=*, itemsep=0.1in, label={\color{brandred}\textbullet}]
        \item Relied on a random train/test split.
        \item Used accuracy as the only metric.
        \item Normalized on the full dataset before splitting.
        \item Assumed past data represents future.
    \end{itemize}

\end{minipage}
\hfill
\begin{minipage}[t]{0.45\textwidth}

    {\large\color{brandgreen}\textbf{What I Should Have Done}}

    \begin{itemize}[leftmargin=*, itemsep=0.1in, label={\color{brandgreen}\textbullet}]
        \item Time-based split for temporal data.
        \item Business-aligned metrics (precision/recall, F-beta).
        \item Fit scalers only on training split.
        \item Implement drift detection and periodic retraining.
    \end{itemize}

\end{minipage}

\vspace{0.3in}

\begin{darkbox}

{\large\textbf{The Reality Check}}

\vspace{0.05in}

The model wasn't wrong – it was evaluated on a lie. The validation set leaked information and ignored the real distribution.

\end{darkbox}

\vfill
\newpage
% ==========================================================
% SLIDE 4: HOW IT SHOWED UP (SINGLE PAGE)
% ==========================================================

\section*{How It Showed Up}

\vspace{0.08in}

{\small No crashes. No exceptions. Just three silent failures that made the model confidently wrong.}

\vspace{0.12in}

% --- FAILURE 1: DATA LEAKAGE ---
\begin{infobox}[top=6pt, bottom=6pt, left=12pt, right=12pt]
{\normalsize\color{brandblue}\textbf{1. Data Leakage}} \hfill {\footnotesize\color{brandorange}\textit{The model saw the test.}}

\vspace{0.04in}
{\footnotesize\textbf{Punchline:} Scaling before split $=$ test statistics leaked into training.}

\vspace{0.03in}
{\footnotesize\color{brandgray}\textit{Example: \texttt{scaler.fit(X\_all)} instead of \texttt{scaler.fit(X\_train)} — the model memorized answers instead of learning patterns.}}

\vspace{0.03in}
{\footnotesize\color{brandred}\textbf{Human fix:} Verify every preprocessing step respects the train/test boundary.}
\end{infobox}

\vspace{0.08in}

% --- FAILURE 2: EVALUATION GAP ---
\begin{infobox}[top=6pt, bottom=6pt, left=12pt, right=12pt]
{\normalsize\color{brandblue}\textbf{2. Evaluation Gap}} \hfill {\footnotesize\color{brandorange}\textit{Accuracy $\neq$ truth.}}

\vspace{0.04in}
{\footnotesize\textbf{Punchline:} 98\% accuracy meant nothing when the model missed 95\% of the rare, high-value class.}

\vspace{0.03in}
{\footnotesize\color{brandgray}\textit{Example: 0\% recall on fraud detection — the model just predicted ``not fraud'' every time.}}

\vspace{0.03in}
{\footnotesize\color{brandred}\textbf{Human fix:} Let domain experts define what ``good'' means before training begins.}
\end{infobox}

\vspace{0.08in}

% --- FAILURE 3: SILENT DRIFT ---
\begin{infobox}[top=6pt, bottom=6pt, left=12pt, right=12pt]
{\normalsize\color{brandblue}\textbf{3. Silent Drift}} \hfill {\footnotesize\color{brandorange}\textit{Yesterday's data $\neq$ today's world.}}

\vspace{0.04in}
{\footnotesize\textbf{Punchline:} User behavior changed, but the validation set stayed frozen in the past.}

\vspace{0.03in}
{\footnotesize\color{brandgray}\textit{Example: Pre-pandemic shopping habits became irrelevant after lockdowns — but the model kept serving stale recommendations.}}

\vspace{0.03in}
{\footnotesize\color{brandred}\textbf{Human fix:} Monitor for drift and decide when the world has changed enough to retrain.}
\end{infobox}

\vspace{0.1in}

% --- CLOSING KICKER ---
\begin{darkbox}[top=6pt, bottom=6pt, left=12pt, right=12pt]
{\normalsize\textbf{The Real Danger}}

\vspace{0.03in}
{\footnotesize AI doesn't crash — it convinces you it's right. Silent failures prove why AI needs human guardrails at every stage.}
\end{darkbox}

\vfill
\newpage


% ==========================================================
% SLIDE 5: ROOT CAUSE ANALYSIS
% ==========================================================
\section*{Root Cause Analysis}

\vspace{0.05in}

{\fontsize{11}{13}\selectfont\color{brandgray}
The failure was a compounding of technical shortcuts and missing process controls.
}

\vspace{0.2in}

\begin{minipage}[t]{0.47\textwidth}

{\large\color{brandorange}\textbf{Technical Causes}}

\vspace{0.1in}
\begin{itemize}[leftmargin=*, itemsep=0.15in, label={\color{brandblue}\textbullet}]
    \item \textbf{Preprocessing before split} $\rightarrow$ Data leakage (Slide 4.1)
    \item \textbf{Metric misalignment} $\rightarrow$ Evaluation gap (Slide 4.2)
    \item \textbf{No temporal awareness} $\rightarrow$ Silent drift (Slide 4.3)
    \item \textbf{No monitoring} $\rightarrow$ All three undetected
\end{itemize}

\end{minipage}
\hfill
\begin{minipage}[t]{0.47\textwidth}

{\large\color{brandorange}\textbf{Process Causes}}

\vspace{0.1in}
\begin{itemize}[leftmargin=*, itemsep=0.15in, label={\color{brandblue}\textbullet}]
    \item \textbf{No domain expert review} – metrics didn't reflect business priorities.
    \item \textbf{Validation not simulating production} – missing real data distribution.
    \item \textbf{No A/B testing} – no safe rollout to catch failures.
    \item \textbf{No feedback loop} – user complaints never reached the model team.
\end{itemize}

\end{minipage}

\vfill
\newpage

% ==========================================================
% SLIDE 6: WHAT I CHANGED
% ==========================================================

\section*{What I Changed}

\vspace{0.1in}

Three practical changes turned validation from a vanity metric into a trustworthy gate.

\vspace{0.15in}

\begin{darkbox}
\textbf{1. Leakage-proof preprocessing}

Fit scalers, encoders, and imputers \textit{only} on the training split. Use \texttt{sklearn.pipeline.Pipeline} to bundle all steps.
\end{darkbox}

\vspace{0.15in}

\begin{darkbox}
\textbf{2. Business-aligned metrics}

Replace accuracy with precision, recall, F-beta, or expected value. Choose the metric that reflects the cost of being wrong.
\end{darkbox}

\vspace{0.15in}

\begin{darkbox}
\textbf{3. Continuous validation \& drift detection}

Set up periodic evaluation on fresh data, monitor feature distributions, and retrain when drift exceeds a threshold.
\end{darkbox}

\vfill
\newpage

% ==========================================================
% SLIDE 7: THE IMPLEMENTATION
% ==========================================================

\section*{The Implementation}

\vspace{0.1in}

These snippets show the core changes that prevent silent evaluation failures.

\vspace{0.15in}

\textbf{1. Leakage-proof pipeline (Python)}

\begin{codebox}{Python}
# pipeline.py
from sklearn.pipeline import Pipeline
from sklearn.preprocessing import StandardScaler
from sklearn.ensemble import RandomForestClassifier

pipeline = Pipeline([
    ('scaler', StandardScaler()),
    ('model', RandomForestClassifier())
])

# Fit only on train
pipeline.fit(X_train, y_train)

# Transform test without re-fitting
y_pred = pipeline.predict(X_test)
\end{codebox}

\vspace{0.15in}

\textbf{2. Metric selection \& drift check (Python)}

\begin{codebox}{Python}
# metrics.py
from sklearn.metrics import classification_report, fbeta_score
print(classification_report(y_true, y_pred))

# Custom scorer: recall twice as important
business_score = fbeta_score(y_true, y_pred, beta=2)
print(f"F2 Score: {business_score:.3f}")

# Drift check (Population Stability Index)
psi = compute_psi(expected_dist, actual_dist)
if psi > 0.25:
    trigger_retraining()
\end{codebox}

\vfill
\newpage

% ==========================================================
% SLIDE 8: THE BIG TAKEAWAY
% ==========================================================

\vspace*{0.35in}

\begin{center}

{\fontsize{14}{16}\selectfont\color{brandblue}\textbf{THE MLOPS TAKEAWAY}}

\vspace{0.3in}

\begin{minipage}{0.88\textwidth}
\centering

{\color{brandblue}\rule{6cm}{1pt}}

\vspace{0.2in}

{\fontsize{24}{28}\selectfont\bfseries\color{brandnavy}
A great validation score means\\[0.1in]
nothing if your validation set lies.
}

\vspace{0.15in}

{\fontsize{12}{14}\selectfont\color{brandgray}
Leakage, metric misalignment, and drift silently corrupt your evaluation. If your validation doesn't mirror production, you're optimizing for a fantasy.

\vspace{0.1in}
Make validation a simulation of reality – not a formality.
}

\vspace{0.2in}

{\color{brandblue}\rule{6cm}{1pt}}

\vspace{0.3in}

{\fontsize{14}{16}\selectfont\bfseries\color{brandnavy}
What's the biggest evaluation surprise\\
you've caught in production?
}

\end{minipage}

\end{center}

\vfill

% Full footer
\begin{center}

\rule{0.6\textwidth}{0.5pt}

\vspace{0.12in}

{\fontsize{11}{13}\selectfont\color{brandgray}
Building AI systems by understanding the mathematics,\\
engineering, and intuition behind every architecture.
}

\vspace{0.12in}

{\fontsize{10}{12}\selectfont\color{brandgray}
\begin{tabular}{@{}l l}
GitHub: & \texttt{github.com/MrHeaven1y} \\
Portfolio: & \texttt{mrheaven1y.github.io} \\
LinkedIn: & \texttt{linkedin.com/in/dibayendu-mukherjee-bb897b267}
\end{tabular}
}

\vspace{0.12in}

{\fontsize{14}{16}\selectfont\color{brandnavy}\textbf{Dibayendu Mukherjee}}

\end{center}

\vspace*{0.2in}


% ==========================================================
% SLIDE 9: COMING UP NEXT
% ==========================================================

\thispagestyle{empty}
\vspace*{0.5in}

\begin{center}
    {\fontsize{14}{16}\selectfont\color{brandblue}\textbf{COMING UP NEXT...}}
    
    \vspace{0.4in}
    
    {\fontsize{32}{38}\selectfont\color{brandnavy}\bfseries
    What Actually Breaks When\\[0.15in]
    You Scale Training Across\\[0.15in]
    Multiple GPUs
    }
    
    \vspace{0.4in}
    
    \begin{minipage}{0.7\textwidth}
        \centering
        \color{brandgray}\fontsize{14}{18}\selectfont
        Distributed training, memory bottlenecks,\\
        and the hidden cost of parallelism.
    \end{minipage}
\end{center}

\vfill

\begin{center}
    \rule{0.6\textwidth}{0.5pt}
    
    \vspace{0.15in}
    
    {\fontsize{11}{13}\selectfont\color{brandgray}
    More production-grade MLOps deep dives coming soon.
    }
    
    \vspace{0.12in}
    
    {\fontsize{14}{16}\selectfont\color{brandnavy}\textbf{Dibayendu Mukherjee}}
\end{center}

\vspace*{0.2in}

\newpage

\end{document}